\documentclass[12pt]{article}
\usepackage[english]{babel}
\usepackage{graphicx}
\usepackage{framed}
\usepackage[normalem]{ulem}
\usepackage{amsmath}
\usepackage{amsthm}
\usepackage{amssymb}
\usepackage{amsfonts}
\usepackage{enumerate}
\usepackage{graphicx} % Allows including image
\usepackage{hyperref}
\hypersetup{
  linkcolor  = blue,
  citecolor  = blue,
  urlcolor   = blue,
  colorlinks = true,
} % color setup
\usepackage[utf8]{inputenc}
\usepackage[margin=1in]{geometry}
\usepackage{float} % for "H" placement specifier
\usepackage[authoryear]{natbib}
\usepackage{indentfirst}
\usepackage[dvipsnames]{xcolor}
\usepackage{float}
\usepackage{array}


\usepackage[T1]{fontenc}
\usepackage{tikz}
\newcommand{\fatterdot}{\raisebox{0.15ex}{\tikz\filldraw[black,x=1.5pt,y=1.5pt] (0,0) circle (1);}}

\DeclareUnicodeCharacter{FB01}{fi}
\DeclareUnicodeCharacter{FF0C}{,}
\DeclareMathOperator*{\Prob}{P}
\DeclareMathOperator*{\argmax}{arg\,max}
\DeclareMathOperator*{\argmin}{arg\,min}
\setlength\extrarowheight{1pt}

\newcommand\red[1]{{\color{red}#1}}

\bibliographystyle{plainnat}

\newtheorem{theorem}{Theorem}[section]
\newtheorem{corollary}{Corollary}[theorem]
\newtheorem{lemma}[theorem]{Lemma}
\newtheorem{proposition}[theorem]{Proposition}


\title{Foundations of Full House Methodology}
\author{}

\date{}

\begin{document}

\maketitle


We now motivate the structure of the Full House Model. Let $N_i$ denote the size of population $i$. We will suppose the every individual $k = 1,\ldots, N_i$ has an underlying talent value $X_{i,1}, \ldots, X_{i,N_i} \stackrel{\text { iid }}{\sim} F_{X}$. We denote $X_{i,(j)}$ as the $j$th ordered latent aptitude in year $i$.
We define $g_{i}\left(\cdot, \cdot\right)$ as the system inclusion function, where $g_{i}\left(X_{i,j}, \textbf{X}_{i,-j}\right)=1$ indicates that individual $j$ is included in the $i$th system, and $g_{i}\left(X_{i,j}, \textbf{X}_{i,-j}\right)=0$ indicates that individual $j$ is not included in the $i$th system. Let $\textbf{X}_{i,-j}$ be the vector of all latent aptitudes not including the component $j$. 

%We will assume that the most talented individuals in year $i$ are active players in the MLB so that 
%\begin{equation} \label{eq:MLBselection}
%	g_{i}\left(X_{i, j}, \mathbf{X}_{i,-j}\right) = 1\left(X_{i,j} \geq X_{i,(N_i-n_i + 1)}\right),
%\end{equation}
%where $n_i$ is the number of active MLB players in year $i$ and $1(\cdot)$ is the indicator function.

%A general formulation for $g$ is given by,
%\begin{equation} \label{eq:selection}
%	g_{i}\left(X_{i, j}, \mathbf{X}_{i,-j}\right)=1\left( X_{i,j} \in { l_{i,1}(X_{i, j},\mathbf{X}_{i,-j}), l_{i,2}(X_{i, j},\mathbf{X}_{i,-j}), \dots, l_{i,n_i}(X_{i, j},\mathbf{X}_{i,-j})  }  \right),
%\end{equation}
%where $l_{i,k}(X_{i, j},\mathbf{X}_{i,-j}) = X_{i, (r(N_i, n_i, k))}, k = 1, 2, \dots, n_i$, $n_i$ is the number of individuals included in the $i$th system, $r(N_i, n_i, k)$ is the index function that selects components. 

%We can choose any index functions we need to select the desirable components. For example, suppose $r(N_i, n_i, k) = N_i - n_i + k$, $k = 1, 2, \dots, n_i$, then the inclusion function $g_i$ will select the top $n_i$ components from the population $N_i$ into the $i$th system. 

We now suppose that an individual $j$ included in system $i$ has an observed statistics $Y_{i,j}$. We suppose that $Y_{i,1}, \ldots, Y_{i,n_i} \overset{iid}{\sim} F_{Y_i}$ where $F_{Y_i}$ will be continuous, and $n_i < N_i$ is the number of individuals included in system $i$. We denote $Y_{i,(j)}$ as the $j$th ordered statistic for system $i$.

The key to the Full House Model is connecting talent values $X$ with observed statistics $Y$. To do this we assume that we can form pairs $(Y_{i,(j)}, X_{i,(r(N_i, n_i, j))})$, $j = 1,\ldots,n_i$, where $r(N_i, n_i, k)$ is the index function that corresponds to the selected components.  %\textcolor{red}{Needs to be more general.} %As an example, the highest performer in the MLB in year $i$ as judged by values $Y_{i,j}$ will be assumed to have the highest talent score in year $i$.

\bigskip
We now demonstrate how the Full House Model works in a parametric setting. Consider the pair $(Y_{i,(j)}, X_{i,(r(N_i, n_i, j))})$ 
%(\textcolor{red}{Needs to be more general, this currently corresponds to the case where all the highest latent aptitudes are in the system.})
and suppose that the distribution corresponding to $Y_{i, j}$ from the $i$th system is known to belong to a continuous parametric family indexed by unknown parameter $\theta_{i},$ and let $F_{Y_i}(\cdot \mid \theta_i)$ be a parametric CDF with parameter $\theta_i \in \mathbb{R}^{p_{Y_i}}$. We can estimate $\theta_i$ with $\hat{\theta_i}$ and plug the estimator into the CDF $F_{Y}(\cdot \mid \hat{\theta_i})$.

In order to connect the  $Y_{i,(j)}$ and $X_{i,(r(N_i, n_i, j))}$ and obtain an estimate of the underlying talents,  we will make use of the following classical order statistics properties, 
$$
\begin{aligned}
F_{Y_i}\left(Y_{i,(j)} \mid \theta_i\right) & \sim U_{i,(j)}, & & F_{Y_i}\left(Y_{i,(j)} \mid \hat{\theta_i}\right) \approx U_{i,(j)}, \\
F_{Y_{i,(j)}}\left(Y_{i,(j)} \mid \theta_i\right) & \sim U_{i,j}, & & F_{Y_{i,(j)}}\left(Y_{i,(j)} \mid \hat{\theta_i}\right) \approx U_{i,j},
\end{aligned}
$$
where $\approx$ means approximately distributed, $\sim$ means distributed as, $U_{i,j} \sim U(0,1)$, and $U_{i,(j)} \sim \operatorname{Beta}(j, n_i+1-j)$ and the quality of the approximation in the right hand side depends upon the estimator $\hat{\theta}_i$ and the sample size. 

We now connect the order statistics to the underlying talent distribution that comes from a population with $N_i \geq n_i$ observations when $F_{X}$ is known. This connection is established with 
\begin{equation} \label{eq:trueparX}
F_{X_{i, (r(N_i, n_i, j))}}^{-1}\left(F_{U_{i,(j)}}\left(F_{Y_{i}}\left(Y_{i,(j)} \mid \theta_i\right)\right)\right) \sim F_{X_{i, (r(N_i, n_i, j))}}^{-1}\left(F_{U_{i,(j)}}\left(U_{i,(j)}\right)\right) = X_{i, (r(N_i, n_i, j))}.
\end{equation}
We estimate the above with
%$$
%F_{X_{i, (r(N_i, n_i, j))}}^{-1}\left(F_{U_{i,(j)}}\left(F_{Y_{i}}\left(Y_{i,(j)} \mid \hat{\theta_i}\right)\right)\right) \approx F_{X_{i, (r(N_i, n_i, j))}}^{-1}\left(F_{U_{i,(j)}}\left(U_{i,(j)}\right)\right) = X_{i, (r(N_i, n_i, j))}.
%$$
%When $r(N_i, n_i, j) = N_i - n_i + j$, $j = 1, \dots, n_i$, as in our application where the most talented baseball players are assumed to be in the MLB, then this relation becomes
\begin{equation} \label{eq:parestX}
	F_{X_{i, (r(N_i, n_i, j))}}^{-1}\left(F_{U_{i,(j)}}\left(F_{Y_{i}}\left(Y_{i,(j)} \mid \hat{\theta}_i\right)\right)\right) \approx F_{X_{i, (r(N_i, n_i, j))}}^{-1}\left(F_{U_{i,(j)}}\left(U_{i,(j)}\right)\right) = \widehat{X}_{i, (r(N_i, n_i, j))}.
\end{equation}
Similarly, we can obtain $Y$ from $X$. First note that $F_X\left(X_{i, (r(N_i, n_i, j))}\right) \sim \text{Beta}(r(N_i, n_i, j), N_i +1 - r(N_i, n_i, j))$ for $j = 1,\ldots,n_i$ where, for simplicity, we we denote $V(i,j) \overset{d}{=} \text{Beta}(r(N_i, n_i, j), N_i +1 - r(N_i, n_i, j))$. Now, 
\begin{equation} \label{eq:trueparY}
 F_{Y_{i, (j)}}^{-1}\left( F_{V(i,j)}\left(F_X\left(X_{i, (r(N_i, n_i, j))}\right)\right) | \theta_i\right) = Y_{i, (j)},
\end{equation}
where we estimate the above with $F_{Y_{i, (j)}}^{-1}\left( F_{V(i,j)}\left(F_X\left(\widehat{X}_{i, (r(N_i, n_i, j))}\right)\right) | \hat{\theta}_i\right) = \widehat{Y}_{i, (j)}$ for $j = 1,\ldots,n_i$.


\bigskip\bigskip
\textcolor{red}{Not sure where cross-contextual comparisons are needed.} We can now reverse engineer the process above to obtain estimated statistics in another context.
Consider the the pair $(Y_{i,(j)}, X_{i,(r(N_i, n_i, j))})$ in the $i$th population. We first put the aptitude value $X_{i,(r(N_i, n_i, j))}$ obtained by \eqref{eq:parestX} in a new population $m$ and reverse the process to obtain the hypothetical statistics in system $m$, which we will denote as $Y_{i,j,m}$. %The distribution $F_X$ is known. 
More formally, $Y_{i,j,m}$ are computed as follows when baseball statistics are estimated parametrically,
\begin{equation} \label{eq:parrev}
	\widehat{Y}_{i,j,m} = F_{Y_{m}}^{-1} \left( F^{-1}_{U_{m, (l_{i,j,m})}}
\left(F_{X_{m, (r(N_m, n_m, l_{i,j,m}))}}\left(\widehat{X}_{i, (r(N_i, n_i, j))}  \right) \right)| \hat{\theta}_m\right),
\end{equation}
where $l_{i,j,m}$ is the rank of $\widehat{X}_{i,(r(N_i, n_i, j))}$ among the values
$$
\left\{\widehat{X}_{i,(r(N_i, n_i, j))}, \widehat{X}_{m,(r(N_m, n_m, t))}: t=1, \ldots, n_{m}\right\}.
$$
The estimates \eqref{eq:parrev} correspond to 
\begin{equation} \label{eq:truparY}
	Y_{i,j,m} = F_{Y_{m}}^{-1} \left( F^{-1}_{U_{m, (l_{i,j,m})}}
\left(F_{X_{m, (r(N_m, n_m, l_{i,j,m}))}}\left(X_{i, (r(N_i, n_i, j))} \right) \right)\right),
\end{equation}
where $l_{i,j,m}$ is the rank of $X_{i,(r(N_i, n_i, j))}$ among the values
$$
\left\{X_{i,(r(N_i, n_i, j))}, X_{m,(r(N_m, n_m, t))}: t=1, \ldots, n_{m}\right\},
$$
and $X_{i,(r(N_i, n_i, j))}$ is computed via \eqref{eq:trueparX}.




\section*{Paradoxes}


Consider two populations $i = 1,2$ with sizes $N_1$ and $N_2$, and systems of sizes $n_1$ and $n_2$. We define a paradox (or a Gould paradox) when 
$$
  X_{1,(N_1)} > X_{2,(N_2)} \quad \text{and} \quad Y_{1,(n_1)} < Y_{2,(n_2)}
$$
both occur. A central interest is evaluating the probability of such a paradox, 
\begin{equation} \label{eq:simpleparadox}
  \Prob\left(X_{1,(N_1)} > X_{2,(N_2)}, Y_{1,(n_1)} < Y_{2,(n_2)}\right).
\end{equation}
As a simple illustrative example of some mechanics concerning paradoxes, suppose that $Y_{1,j} < Y_{2,k}$ a.s. Then \eqref{eq:simpleparadox} is equivalent to $\Prob\left(X_{1,(N_1)} > X_{2,(N_2)}\right)$, and we have $\Prob\left(X_{1,(N_1)} > X_{2,(N_2)}\right) \to 1$ as $N_1 \to \infty$.





Consider the setting where $g_{i}\left(X_{i, j}, \mathbf{X}_{i,-j}\right) = 1\left(X_{i,j} \geq X_{i,(N_i-n_i + 1)}\right)$. Then \eqref{eq:trueparY} gives 
$$
  F_{Y_{i, (n_i)}}^{-1}\left( F_{V(i,n_i)}\left(F_X\left(X_{i, (N_i)}\right)\right) | \theta_i\right) = Y_{i, (n_i)} \quad \text{for} \; i = 1,2.
$$
Let $f_i(\fatterdot) = F_{Y_{i, (n_i)}}^{-1}\left( F_{V(i,n_i)}\left(F_X\left(\fatterdot\right)\right) | \theta_i\right)$ and let $g_i(\fatterdot) = f_i^{-1}(\fatterdot)$.
%\begin{align*}
%  f_i(\fatterdot) &= F_{Y_{i, (n_i)}}^{-1}\left( F_{V(i,n_i)}\left(F_X\left(\fatterdot\right)\right) | \theta_i\right),	\\
%  g_i(\fatterdot) &= F_X^{-1}\left(F_{V(i,n_i)}^{-1}\left(F_{Y_{i,(n_i)}}\left(\fatterdot|\theta_i\right) \right) \right),
%\end{align*}
This construction yields,
\begin{align*} 
  \Prob\left(X_{1,(N_1)} > X_{2,(N_2)}, Y_{1,(n_1)} < Y_{2,(n_2)}\right) 
  = \Prob\left(X_{1,(N_1)} > X_{2,(N_2)}, f_1\left(X_{1,(N_1)}\right) < f_2\left(X_{2,(N_2)}\right)\right).
\end{align*}
Let $h_{i,k}(\fatterdot) = (g_i \circ f_k)(\fatterdot)$ when such an operation is defined. Then, 
\begin{align*} 
  &\Prob\left(X_{1,(N_1)} > X_{2,(N_2)}, Y_{1,(n_1)} < Y_{2,(n_2)}\right) 
  = \Prob\left(X_{1,(N_1)} > X_{2,(N_2)}, X_{1,(N_1)} < h_{1,2}\left(X_{2,(N_2)}\right)\right) \\
  &\qquad= \Prob\left( X_{2,(N_2)} < X_{1,(N_1)}  < h_{1,2}\left(X_{2,(N_2)}\right)\right).
\end{align*}


\end{document}


